\documentclass[10pt,a4paper]{amsart}
\usepackage[margin=19mm]{geometry}
\usepackage{amsmath,amssymb,mathtools}
\usepackage[scr=rsfs,cal=euler]{mathalfa}
\usepackage{xfrac}
\usepackage[unicode,hidelinks,bookmarks=false]{hyperref}
\newcommand{\ie}{i.e.\ }
\newcommand{\cf}{cf.\ }
\newcommand{\resp}{respectively\ }
\newcommand{\Subsection}{\subsection}
\providecommand{\nobreakdash}{\nobreak\hbox{-}}
\setlength{\emergencystretch}{2em}
\multlinegap0pt
\usepackage{dsfont}
\usepackage{amssymb}
\usepackage{xcolor}
\usepackage{algorithm}\usepackage{algpseudocode}
\usepackage{bm}
\newtheorem{theorem}{Theorem}
\usepackage{cleveref}
\crefname{theorem}{Theorem}{Theorems}
\renewcommand{\Im}[1]{\mathrm{Im}(#1)}
\renewcommand{\Re}[1]{\mathrm{Re}(#1)}
\newcommand{\new}[1]{{\textcolor{black}{#1}}} 
\title{Epstein zeta method for many-body lattice sums}
\author{Andreas A. Buchheit, Jonathan K. Busse}
\date{Source: arXiv:2504.11989v2 (CC BY 4.0)}
\begin{document}
\maketitle

\noindent\emph{Source and licence.} Epstein zeta method for many-body lattice sums. Version arXiv:2504.11989v2 (CC BY 4.0), distributed by arXiv under the Creative Commons Attribution 4.0 International licence, http://creativecommons.org/licenses/by/4.0/, as recorded in the arXiv licence metadata for this exact version and shown on the abstract page https://arxiv.org/abs/2504.11989v2. Full licence notice: \texttt{licenses/arxiv-2504.11989-CC-BY-4.0.txt}.\par\smallskip

\noindent\textit{Editorial scope.} Counted unit: the article's theorem holreg (source0.tex lines 362--374). The article's complete original proof of it is delivered with it (source0.tex lines 392--480, one \texttt{proof} environment of 89 lines, which the article places away from the statement). No mathematics is written, repaired or re-derived: the statement and the proof are the article's own lines, in the article's own notation, and no retained line is substituted or re-flowed.  No further source-local context is needed: every cross-reference of every retained line resolves inside the two delivered spans.  Nothing was omitted from any retained span, so the package carries no omission record.  No retained line contains a citation, so the wrapper carries no bibliography and no bibliography entry.  No retained line contains a figure environment, an image-inclusion command or drawing code, so no image is delivered and the package declares no asset dependency.  Editorial formatting only, in addition: an AMS \texttt{amsart} preamble that restates the article's own declarations and macros used by the retained lines, namely the environments \texttt{theorem} on separate counters, together with the standard packages those lines need.  The retained environments print in this document as Theorem 1 in order of appearance, whereas the article prints the same items with the numbering of its own full document.  The complete native source of the article as distributed in the arXiv e-print for this exact version is bundled unchanged as \texttt{sources/176638/source0.tex}.
\begin{theorem}
\label{holreg}
Let $\Lambda$ be a $d$-dimensional lattice and $\nu\in\mathds C$. 
Denote by 
$$
D_L=\{\bm q\in\mathds C^d:|\Re{\bm q}-\bm z|>|\Im{\bm q}|\ \forall \bm z\in L\}.
$$
the intersection of $d$-dimensional complex cones with origins at $L\subseteq \mathds R^d$.
Then, $Z_{\Lambda,\nu}(\bm \cdot)$ can be holomorphically extended to $D_{\Lambda^\ast}$
and $Z^{\rm reg}_{\Lambda,\nu}(\bm \cdot)$ 
can be holomorphically extended to 
$D_{\Lambda^*\setminus\{\bm 0\}}$.
\end{theorem}

\begin{proof}
\new{Recall that \(f\) is a product of Epstein zeta functions and that
\[
\bm w(\bm v)
=
A^{-T}\bigl(\sigma^j(p_1v_1,\dots,p_{d-1}v_{d-1},1)^T\bigr),
\]
where \(\sigma\) denotes a cyclic permutation and \(p_i\in\{\pm1\}\). We begin with the holomorphic extension of \(f\) itself and then incorporate, step by step, the permutation, sign changes, rescaling, and dimension reduction.}

\new{The Epstein zeta function $Z_{\Lambda,\nu_i}$ extends holomorphically to $D_{\Lambda^*}$ by \Cref{holreg}. Thus, $f$ extends to a holomorphic function on $D_{\Lambda^*}$
as the product of holomorphic functions and,
by the composition theorem,
$f(A^{-T} \bm q)$ is holomorphic on 
$\bm q\in 
A^TD_{\Lambda^*}$. As $\Lambda^\ast=A^{-T} \mathds Z^d$, we find after inserting in \Cref{holreg} that
\[
A^T D_{\Lambda^*}=\{\bm q\in\mathds C^d:|A^{-T}(\Re{\bm q}-\bm z)|>|A^{-T}\Im{\bm q}|\ \forall \bm z\in \mathds Z^d\}.
\]}
\new{Consider the simplified set
\[
K
=
\{\bm q\in\mathds C^d:\ |\Re{\bm q}-\bm z|>\kappa(A)|\Im{\bm q}|
\quad \forall \bm z\in\mathds Z^d\},
\]
which depends only on the condition number \(\kappa(A)\). We show that $K\subset A^T D_{\Lambda^\ast}$.
For \(\bm q\in K\) and \(\bm z\in\mathds Z^d\), the singular-value bounds
\[
|A^{-T}\bm k|\ge \sigma_{\min}(A^{-T})|\bm k|,
\qquad
|A^{-T}\bm k|\le \sigma_{\max}(A^{-T})|\bm k|,
\qquad \bm k\in\mathds R^d,
\]
together with
$
\sigma_{\min}(A^{-T})=1/\sigma_{\max}(A)$, and
$
\sigma_{\max}(A^{-T})=1/\sigma_{\min}(A),
$
yield
\begin{align*}
|A^{-T}(\Re{\bm q}-\bm z)|
\ge \frac{1}{\sigma_{\max}(A)}|\Re{\bm q}-\bm z|
> \frac{\kappa(A)}{\sigma_{\max}(A)}|\Im{\bm q}|
&= \frac{1}{\sigma_{\min}(A)}|\Im{\bm q}|\\
&\ge |A^{-T}\Im{\bm q}|.
\end{align*}
Hence, we obtain \(\bm q\in A^T D_{\Lambda^\ast}\), proving that \(K\subset A^T D_{\Lambda^\ast}\). Therefore, \(f(A^{-T}\bm q)\) is holomorphic on \(\bm q\in K\).
Permutations and sign change can further be included via the signed permutation matrix \[Q_{j,\bm p}=\sigma^j \mathrm{diag}(p_1,\dots,p_{d-1},1),\] with $p_i= \pm 1$. Then $f(A^{-T} Q_{j,\bm p}\bm q)$ is also holomorphic on $K$, as the set inherits invariance under entry-wise sign changes and permutations of vector entries from the relation $Q_{j,\bm p}\mathds Z^d =\mathds Z^d$. Moreover, for $0<u<1/2$, we have that $f(u A^{-T} Q_{j,\bm p}\bm q)$ is holomorphic on $K/u$.}

\new{Finally, we transfer this result to holomorphy of $f(u\bm w(\bm v))$ in $\bm v\in\mathds C^{d-1}$. By the preceding argument, $f(u\bm w(\bm v))$ is holomorphic in $\bm v$ on
\[
\left\{
\bm v\in\mathds C^{d-1}:
\left|\Re{u(\bm v,1)^T}-\bm z\right|
>
\kappa(A)\left|\Im{u(\bm v,1)^T}\right|
\quad \forall \bm z\in\mathds Z^d
\right\}.
\]
Since $u>0$, squaring the inequality and separating the last component shows that the domain is equal to
\[
\left\{
\bm v\in\mathds C^{d-1}:
\left|\Re{\bm v}-\tilde{\bm z}/u\right|^2
+
\left(1-z_d/u\right)^2
>
\kappa(A)^2 |\Im{\bm v}|^2
\quad
\forall\, \tilde{\bm z}\in\mathds Z^{d-1},\ z_d\in\mathds Z
\right\}.
\]
We now have that $\left|\Re{\bm v}-\tilde{\bm z}/{u}\right|^2 \ge 0$,
and, since $0<u<1/2$,
\[
\left(1-z_d/u\right)^2 \ge 1
\qquad \forall\, z_d\in\mathds Z,
\]
where the bound is reached at $z_d=0$.
Therefore, the above domain contains
\[
\Omega=\left\{
\bm v\in\mathds C^{d-1}:
|\Im{\bm v}|<1/\kappa(A)
\right\}.
\]
on which $f(u\bm w(\bm v))$ is therefore holomorphic as stated.}
\end{proof}

\end{document}
